\section*{Chapter \ref{ch:s2:supply-trades} --- Supply Trades}

\begin{solution}{exo:s2:supply-trades:1}
$0.07 \times 40 = 2.8$ basis points.
\end{solution}

\begin{solution}{exo:s2:supply-trades:2}
$3.44/17 \times \sqrt{12} = 0.70$.
\end{solution}

\begin{solution}{exo:s2:supply-trades:3}
$0.0009 \times \$40$ billion $= \$36$ million to $0.0018 \times \$40$ billion $= \$72$ million.
\end{solution}

\begin{solution}{exo:s2:supply-trades:4}
Knowing the supply is not the same as having the balance sheet to hold it: dealers and investors must make room for the new bonds, and those with
capacity charge for it. End-investors come slowly, so the price moves until they do.
\end{solution}

\begin{solution}{exo:s2:supply-trades:5}
Dealers' balance sheet grew dearer after the crisis, so they take smaller positions and lay them off faster; and investment funds now bid more at
auctions directly, sharing the supply that dealers once absorbed.
\end{solution}

\begin{solution}{exo:s2:supply-trades:6}
The auction closes at 1 p.m., and the end-of-day yield already includes part of the recovery after it; the day's pre-auction cheapening is
partly hidden.
\end{solution}

\begin{solution}{exo:s2:supply-trades:7}
The leg before is unchanged (1.71 basis points per auction); the leg after earns only its costs ($-0.09$); the two together 1.62. The concession paid
before the auction is what dealers are compensated with; the recovery after is a separate bet that end-investors will buy.
\end{solution}

\begin{solution}{exo:s2:supply-trades:8}
Single auctions have a standard deviation of about 17 basis points against a mean of 3.44; three losses of 20 are well within the noise. It takes
dozens of auctions to detect the effect, and as many to detect its disappearance.
\end{solution}

\begin{solution}{pb:s2:supply-trades:1}
\begin{enumerate}
\item The cheapening before a new issue, recovered after; the premium of the newest issue at a maturity.
\item Prices fall before auctions and recover after; a hidden cost of 9 to 18 basis points of size.
\item Issuance drives dealer inventories; dealers are paid by later appreciation; the pay has fallen with higher \omterm{def:s2:swap-spread-and-asset-swap-trades:bscost}{balance-sheet costs} and more fund
participation.
\item They trade at a premium for liquidity and repo value.
\item Pooled, 1.47 basis points up on auction day from ten days before, 0.40 five days after.
\item The 10- and 30-year cheapen before; the 5- and 7-year richen after; the 2-year shows nothing.
\item The rise before grows with size (1.48, 1.69, 2.95 basis points by third); the fall after does not.
\item The part of the auction day before 1 p.m.
\item Short known supply before its sale, long after.
\item 0.07 basis points per billion at average capacity, 70\% given back after.
\item 3.44 per auction, $t$ = 3.1, Sharpe ratio 0.70; each leg about 1.7 and 0.50.
\item Daily noise of 5.5 basis points over ten days swamps a concession of 2.6.
\item \textbf{Named result.} The auction trade earns 3.44 basis points of yield per auction ($t$ = 3.1, Sharpe ratio 0.70 a year), and each billion of
size adds 0.24 basis points on average, though any single auction's result is noise.
\item As much as size: the concession is size over capacity.
\item Around the calendar of large deals, hedged with Treasuries or swaps, taking allocations at the concession.
\item Buying duration before the index's month-end extension, for index trackers' demand.
\item Yields at the trade times, costs per leg, many auctions.
\item The on-the-run roll.
\item Chapter~10's fitted curve excludes the issues this chapter trades around.
\item The issuer, through the concession it pays for quick distribution.
\end{enumerate}
\end{solution}

\begin{solution}{iq:s2:supply-trades:1}
Dealers must absorb the new supply with limited balance sheet, and end-investors arrive slowly; the market prices in the cost of holding the supply
until it is distributed.
\end{solution}

\begin{solution}{iq:s2:supply-trades:2}
Measure yield changes of the auctioned security and close substitutes over windows before and after each auction, relative to a benchmark or a
fitted curve; control for macro news, month-end, other auctions nearby and the level of rates.
\end{solution}

\begin{solution}{iq:s2:supply-trades:3}
Short futures or the when-issued security before the auction in proportion to the expected award, and adjust after the result; plan the
distribution of the bonds to clients.
\end{solution}

\begin{solution}{iq:s2:supply-trades:4}
Macro news around auction dates dominates any one auction; the effect can shrink as more desks trade it; costs and crowding in futures near the
auction.
\end{solution}

\begin{solution}{iq:s2:supply-trades:5}
Announcements (date, size, security), results (high yield, bid-to-cover, allotments) and when-issued quotes, with timestamps; checks on reopenings
and changes to the calendar; a store keyed by CUSIP and auction date.
\end{solution}

\begin{solution}{iq:s2:supply-trades:6}
Over a window of $k$ days the noise is $\sigma\sqrt k$; with $n$ auctions the mean's standard error is $\sigma\sqrt{k/n}$, so $t = c\sqrt n/(\sigma\sqrt
k) = 2$ gives $n = 4k\sigma^2/c^2$: with $\sigma = 5.5$, $k = 5$ and $c = 2.6$, about 90 auctions for the leg before.
\end{solution}
